% !TEX program = xelatex
% ============================================================
%  ندوة عبر الإنترنت — Webinar Template
%  قالب شرائح لندوة إلكترونية مع جدول أعمال ونقاط رئيسية وروابط
%  Compile with: xelatex main.tex (run twice)
% ============================================================
%  Features:
%    - 16:9 widescreen, clean default beamer theme with accent color
%    - Arabic RTL via polyglossia + Amiri font
%    - Speaker intro, agenda, content sections, key takeaways,
%      resources, Q&A, chat-reminder, thank-you slides
%    - English terms wrapped with \textEN{...}
% ============================================================

\documentclass[aspectratio=169,11pt]{beamer}

% ---- Arabic / RTL language support ----
\usepackage{polyglossia}
\setdefaultlanguage[calendar=gregorian,hijricorrection=1]{arabic}
\setotherlanguage[variant=british]{english}

% ---- Fonts (beamer uses sans by default, so set sans to Amiri too) ----
\usepackage[no-math]{fontspec}
\setmainfont[Scale=1]{NotoKufiArabic-Regular.ttf}
\setsansfont[Scale=1]{NotoKufiArabic-Regular.ttf}
\newfontfamily\arabicfont[Script=Arabic,Scale=0.95]{NotoKufiArabic-Regular.ttf}
\newfontfamily\arabicfontsf[Script=Arabic,Scale=0.95]{NotoKufiArabic-Regular.ttf}
\newfontfamily\englishfont[Scale=0.95]{TeX Gyre Termes}
\let\textEN\textenglish

% ---- Core packages ----
\usepackage{graphicx}
\usepackage{booktabs}
\usepackage{tabularx}
\usepackage{tikz}
\usepackage{amsmath}
\usepackage{amssymb}

% ---- Theme & appearance ----
\usetheme{default}
\setbeamertemplate{navigation symbols}{}
\setbeamersize{text margin left=10mm, text margin right=10mm}

% ---- Accent color ----
\definecolor{accent}{HTML}{4527A0}
\setbeamercolor{frametitle}{fg=accent}
\setbeamercolor{title}{fg=accent}
\setbeamercolor{alerted text}{fg=accent}
\setbeamercolor{structure}{fg=accent}
\setbeamertemplate{frametitle}[default][center]

% ---- Metadata ----
% TODO: استبدل بموضوع الندوة وبيانات المقدّم الفعلية
\title{موضوع الندوة}
\subtitle{عنوان فرعي يوضح المحور}
\author{اسم المقدّم}
\institute{الجهة المنظِّمة}
\date{التاريخ — منصة \textEN{[platform]}}

\begin{document}

% ============================================================
% TITLE SLIDE
% ============================================================
\begin{frame}[plain]
\titlepage
\end{frame}

% ============================================================
% ABOUT THE SPEAKER
% ============================================================
\begin{frame}{عن المقدّم}
% TODO: اكتب نبذة مختصرة عن المقدّم
\begin{itemize}
  \item المؤهل العلمي والتخصص.
  \item الخبرة العملية والمجال \textEN{(expertise)}.
  \item الأعمال أو المشاريع البارزة.
  \item روابط التواصل (اختياري).
\end{itemize}
\end{frame}

% ============================================================
% AGENDA
% ============================================================
\begin{frame}{جدول الأعمال}
% TODO: عدّل بنود الجدول حسب الندوة
\begin{enumerate}
  \item مقدمة عن الموضوع.
  \item المحور الأول: المفاهيم الأساسية.
  \item المحور الثاني: التطبيقات العملية.
  \item المحور الثالث: التحديات والحلول.
  \item النقاط الرئيسية والمراجع.
  \item الأسئلة والنقاش.
\end{enumerate}
\end{frame}

% ============================================================
% CONTENT SECTION 1
% ============================================================
\section{المحور الأول: المفاهيم الأساسية}
\begin{frame}{المفاهيم الأساسية}
% TODO: اعرض المفاهيم التأسيسية
\begin{itemize}
  \item تعريف المفهوم الرئيسي.
  \item المصطلحات الأساسية \textEN{(key terms)}.
  \item لماذا يهم هذا الموضوع الجمهور المستهدف.
\end{itemize}
\end{frame}

% ============================================================
% CONTENT SECTION 2
% ============================================================
\section{المحور الثاني: التطبيقات العملية}
\begin{frame}{التطبيقات العملية}
% TODO: اذكر أمثلة عملية
\begin{itemize}
  \item مثال تطبيقي أول مع النتيجة.
  \item مثال تطبيقي ثانٍ \textEN{(use case)}.
  \item دراسة حالة مختصرة.
\end{itemize}
\end{frame}

% ============================================================
% CONTENT SECTION 3
% ============================================================
\section{المحور الثالث: التحديات والحلول}
\begin{frame}{التحديات والحلول}
% TODO: ناقش التحديات الشائعة
\begin{itemize}
  \item التحدي الأول وكيفية التعامل معه.
  \item التحدي الثاني \textEN{(challenge)}.
  \item توصيات عملية للجمهور.
\end{itemize}
\end{frame}

% ============================================================
% KEY TAKEAWAYS
% ============================================================
\begin{frame}{النقاط الرئيسية}
% TODO: لخّص أهم ما يجب أن يستذكره الجمهور
\begin{itemize}
  \item النقطة الأولى: الفكرة الجوهرية.
  \item النقطة الثانية: خطوة عملية قابلة للتطبيق.
  \item النقطة الثالثة: مورد للمتابعة.
  \item النقطة الرابعة: سؤال للتأمل.
\end{itemize}
\end{frame}

% ============================================================
% RESOURCES / LINKS
% ============================================================
\begin{frame}{مصادر وروابط}
% TODO: أضف الروابط الفعلية
\begin{itemize}
  \item مرجع أساسي للموضوع: \textEN{https://example.com/resource1}
  \item أداة أو موقع ذو صلة: \textEN{https://example.com/tool}
  \item ورقة بحثية موصى بها: \textEN{https://example.com/paper}
  \item تسجيل الندوة (بعد البث): \textEN{https://example.com/recording}
\end{itemize}
\end{frame}

% ============================================================
% CHAT REMINDER
% ============================================================
\begin{frame}{تذكير: اطرحوا أسئلتكم في الدردشة}
\begin{center}
\vspace{0.8cm}
{\Large نرحب بأسئلتكم في مربع الدردشة \textEN{(chat)} خلال العرض.}\\[0.8cm]
{\normalsize سيتم الرد على الأسئلة في قسم الأسئلة والنقاش في النهاية.}
\end{center}
\end{frame}

% ============================================================
% Q&A
% ============================================================
\begin{frame}{الأسئلة والنقاش}
\begin{center}
\vspace{1cm}
{\Huge\bfseries\color{accent} الأسئلة والنقاش}\\[0.8cm]
{\large اطرحوا أسئلتكم الآن أو عبر الدردشة.}
\end{center}
\end{frame}

% ============================================================
% THANK YOU + CONTACT
% ============================================================
\begin{frame}[plain]
\begin{center}
\vspace{1.2cm}
{\Huge\bfseries\color{accent} شكراً لكم}\\[1cm]
{\Large للمتابعة والتواصل}\\[0.8cm]
{\large اسم المقدّم}\\
{\normalsize البريد: \textEN{email@example.com}}\\
{\normalsize الموقع: \textEN{https://example.com}}
\end{center}
\end{frame}

\end{document}
